\documentclass[10pt,journal,compsoc]{IEEEtran}
\usepackage[utf8]{inputenc}
\usepackage{amsmath,amsfonts,amssymb}
\usepackage{graphicx}
\usepackage{cite}
\usepackage{algorithmic}
\usepackage{textcomp}

\begin{document}

\title{The Apex Architecture for Generative Print: Integrating API Workflows, Tiled SUPIR, and Perceptual Colorimetry}

\IEEEtitleabstractindextext{%
\begin{abstract}
This paper formalizes a rigorous generative poster pipeline that bridges API-driven ideation with local typography and upscaling. We introduce a layered compositing model accounting for physical print constraints, utilizing measure theory on a defined probability space to define text density via cross-attention probability maps. We define a robust system architecture bridging Ideogram's rasterized RGB output with localized text generation, entirely bypassing intermediate VAE re-encoding to preserve aesthetic detail. To ensure commercial print legibility, structural constraints are augmented with post-generation luminance contrast heuristics based on WCAG ratios and perceptual divergence. To address VRAM limits, we implement Tiled Diffusion upscaling with shared attention for semantic structural alignment, alongside Latent Space Gaussian Blending over the intersection of patch interiors to enforce $C^{\infty}$ continuity. High-frequency gradients are constrained via Sobolev norm regularization with $s>2$. Finally, we formalize RGB-to-CMYK conversion as a projection minimizing non-metric perceptual divergence governed by 3D LUT interpolation and UCR/GCR black generation, formulating total ink limits and gamuts as rigorous topological boundary conditions.
\end{abstract}
}

\maketitle
\IEEEdisplaynontitleabstractindextext
\IEEEpeerreviewmaketitle

\newtheorem{theorem}{Theorem}
\newtheorem{lemma}{Lemma}
\newtheorem{definition}{Definition}

\section{Introduction}
The synthesis of high-density print media via generative models bridges continuous latent spaces and discrete physical gamuts. Addressing the dichotomy between scalable API access and local execution, we formalize the 'Apex Architecture'. This multi-agent workflow utilizes Ideogram 2.0 via API for foundational aesthetics, followed by local inpainting with AnyText 2 and ControlNet for precise typography, Tiled SUPIR for memory-efficient upscaling, and robust 3D LUT-based Perceptual Colorimetry.

\section{System Architecture: API Handoff and Resolution Bridging}
Bridging remote APIs and local solvers requires robust pipeline management. Ideogram API generations return rasterized RGB pixels, circumventing the need to maintain a history of latent tensors. Crucially, to prevent degrading the API's aesthetic details, we explicitly avoid VAE-encoding the RGB output into a latent space prior to upscaling, nor do we downsample the Ideogram output to match AnyText's native resolution. Instead, AnyText operates on a separate high-resolution canvas, and its localized text generations are merged directly into the full-resolution API base image using precise alpha masks.

\section{Measure-Theoretic Typographic Compositing and Legibility}
Commercial print demands adherence to physical constraints: bleed $B$, trim $T$, and safe zones $S$. We model the canvas as a layered measure space, utilizing a compositing model $C = \sum_{k} \alpha_k L_k$. To ensure commercial print legibility, these spatial constraints are augmented with a post-generation contrast analysis evaluating WCAG ratios or $\Delta E_{00}$ between generated text and the background.

\begin{definition}
Let $(\Omega, \mathcal{F}, \mathbb{P})$ be a formal probability space where $\Omega = [0,1]^2$ represents spatial coordinates. Let $\mathcal{F}$ be the Borel $\sigma$-algebra, and $\mu$ the Lebesgue measure. Let $T_i \subset S$ be the bounding box for the $i$-th typographic element, and let $k \in \mathcal{V}_{\text{text}}$ represent the corresponding text tokens. The text density measure $\rho(T_i)$ is defined using cross-attention random variables $M_k(x)$ corresponding strictly to text tokens: $\rho(T_i) = \int_{T_i} \sum_{k \in \mathcal{V}_{\text{text}}} \mathbb{E}[M_k(x)] d\mu(x)$, under the normalization constraint that $\int_{T_i} \rho d\mu = 1$.
\end{definition}

\begin{theorem}
Let $A(x, y)$ be the cross-attention map between spatial query $x \in \Omega$ and text token $y \in \mathcal{V}_{\text{text}}$. Introducing proximal operators or subgradients onto the attention scores rigorously constraints the support while preserving the differentiability of the score function. Applying these constraints yields an overwhelming concentration of measure within $T_i$.
\end{theorem}
\textit{Proof Sketch.} Utilizing proximal operators smoothly projects the score function onto the constrained manifold. Through Lipschitz bounds on the drift term and standard concentration inequalities (e.g., Chernoff bounds), one formally derives an exponentially bounded tail outside the target support, ensuring typographical elements strictly reside within the safe zone $S$.

\section{VRAM-Aware Tiled Upscaling and Semantic Seam Blending}
To overcome VRAM limits for 77.7 MP outputs, we employ computationally viable Tiled Diffusion. The global image is partitioned into overlapping patches $X$ and $Y$. 

\begin{definition}
To ensure mathematical continuity and semantic structural alignment across seams, we employ Latent Space Gaussian Blending augmented with shared cross-attention over overlapping context windows. Gaussian feathering applies a smooth partition of unity strictly over the intersection of their interiors, $X^\circ \cap Y^\circ$ (which has positive Lebesgue measure). This enforces $C^\infty$ continuity, seamlessly blending latent features while the shared attention enforces structural continuity.
\end{definition}

To resolve the text-upscaling compounding error, we define a strict text masking bypass pipeline. Text regions are rigorously masked out so SUPIR operates entirely un-filtered on non-text areas, preserving the micro-details it was designed to compute without ruining text structures. To mathematically constrain total variation and gradient explosion, we enforce a Sobolev norm bound $||Y||_{H^s} \le C$. By the Sobolev embedding theorem, to bound the first derivative and embed into $C^1$ for $d=2$, we rigorously specify $s > 1 + d/2 = 2$.

\section{Colorimetry: 3D LUTs and CMYK Projections}
\begin{lemma}
The mapping $\Phi: \mathcal{C}_{RGB} \to \mathcal{C}_{CMYK}$ under Perceptual rendering intent acts as a projection onto the bounded CMYK manifold, determined by standard tetrahedral interpolation of a 3D LUT.
\end{lemma}

The mapping $\Phi$ minimizes the non-metric perceptual divergence $\Delta E_{00}$ between distinct manifolds. Crucially, ICC profile Black Generation (UCR/GCR) manages the neutral axis. The total ink limit and physical gamuts impose rigorous boundary conditions on the destination manifold. Specifically, the boundary is the topological boundary of the intersection:
$ \partial \left( [0,1]^4 \cap \{ \vec{c} \mid \vec{c} \cdot \vec{1} \le \tau \} \right) $
where $\vec{c} = (c,m,y,k)$ and $\tau$ is typically $300\%$. This accurately models all printable faces of the CMYK hypercube and explicitly bounds the minimum printable luminance $L^*_{min}$.

\section{Conclusion}
The revised Apex Architecture provides a robust, physically and mathematically sound pipeline. By combining raster-based API handoffs with localized masked typography, proximal operator regularization, structural shared-attention seam blending over interior intersections, un-filtered SUPIR routing with $s>2$ Sobolev bounds, and strict topological CMYK boundary modeling, we guarantee rigorous commercial print synthesis.

\end{document}
